\chapter{Topología de la recta real}\label{ch:b1:topology}

Los límites remiten una y otra vez al mismo vocabulario geométrico:
puntos «próximos a» un \omterm{def:b1:logic:sets}{conjunto}, \omterm{def:b1:logic:sets}{conjuntos} «sin fugas por la
\omterm{def:b1:topology:closure}{frontera}», \omterm{prop:b1:reals:intervals}{intervalos} de los que las sucesiones no pueden escapar.
Este capítulo fija ese vocabulario --- \omterm{def:b1:topology:open}{conjuntos abiertos} y
\omterm{def:b1:topology:closed}{cerrados}, \omterm{def:b1:topology:closure}{interior} y \omterm{def:b1:topology:closure}{clausura}, \omterm{def:b1:topology:dense}{densidad} --- sobre la recta real,
y demuestra la compacidad de los segmentos en su forma secuencial. Las
mismas nociones, en espacios vectoriales normados, son materia del
segundo año; sobre $\R$ están al alcance y son inmediatamente útiles
para el~\cref{ch:b1:continuity}.

\section{Conjuntos abiertos, conjuntos cerrados}

\begin{definition}[Entorno, conjunto abierto]\label{def:b1:topology:open}
Un \omterm{def:b1:logic:sets}{conjunto} $V \subseteq \R$ es un \emph{entorno}\index{entorno} de
$x \in \R$ cuando contiene un \omterm{prop:b1:reals:intervals}{intervalo} $\intoo{x - r}{x + r}$ para
algún $r > 0$. Un \omterm{def:b1:logic:sets}{conjunto} $U \subseteq \R$ es
\emph{abierto}\index{conjunto abierto} cuando es entorno de cada uno de
sus puntos:
\[
\forall x \in U,\ \exists r > 0, \quad \intoo{x - r}{x + r}
\subseteq U .
\]
\end{definition}

\begin{example}\label{ex:b1:topology:open}
Los \omterm{prop:b1:reals:intervals}{intervalos} \omterm{def:b1:topology:open}{abiertos} son \omterm{def:b1:topology:open}{abiertos}: para $x \in \intoo{a}{b}$,
tómese $r = \min(x - a,\, b - x) > 0$. Las semirrectas
$\intoo{a}{+\infty}$ son \omterm{def:b1:topology:open}{abiertas}; $\R$ y $\emptyset$ son \omterm{def:b1:topology:open}{abiertos}
(el segundo, por vacuidad). $\intcc{0}{1}$ \emph{no} es \omterm{def:b1:topology:open}{abierto}:
ningún \omterm{prop:b1:reals:intervals}{intervalo} en torno a $0$ se queda dentro.
\end{example}

\begin{proposition}[Estabilidad de los abiertos]\label{prop:b1:topology:openstable}
Toda unión de \omterm{def:b1:topology:open}{abiertos} es \omterm{def:b1:topology:open}{abierta}; una intersección \emph{finita} de
\omterm{def:b1:topology:open}{abiertos} es \omterm{def:b1:topology:open}{abierta}. Las intersecciones infinitas pueden fallar:
$\bigcap_{n \geq 1} \intoo{-\frac 1n}{\frac 1n} = \{0\}$, que no es \omterm{def:b1:topology:open}{abierto}.
\end{proposition}

\begin{proof}
Unión: si $x \in \bigcup_i U_i$, entonces $x \in U_{i_0}$ para algún
$i_0$, y el \omterm{prop:b1:reals:intervals}{intervalo} que proporciona $U_{i_0}$ queda dentro de la
unión. Intersección finita: si $x \in U_1 \cap \dots \cap U_k$, tómese
$r = \min(r_1, \dots, r_k) > 0$ de los radios que proporciona cada
$U_j$. Para el contraejemplo: todo \omterm{prop:b1:reals:intervals}{intervalo} en torno a $0$ contiene
algún $\frac{1}{n}$ (Arquímedes) y, por tanto, sale de la intersección.
\end{proof}

\begin{example}[Certificar que un conjunto es abierto con radios explícitos]\label{ex:b1:topology:radii}
¿Es \omterm{def:b1:topology:open}{abierto} $U = \{x \in \R : x^2 > 2\}$? Sí, y el certificado se
puede escribir: $U = \intoo{-\infty}{-\sqrt2} \cup
\intoo{\sqrt2}{+\infty}$, unión de dos semirrectas \omterm{def:b1:topology:open}{abiertas}, luego
\omterm{def:b1:topology:open}{abierto} por la~\cref{prop:b1:topology:openstable}. Alternativamente,
razónese punto por punto: para $x \in U$ con $x > \sqrt 2$, tómese
$r = x - \sqrt2 > 0$: todo $y \in \intoo{x - r}{x + r}$ cumple
$y > \sqrt 2$ y, por tanto, $y^2 > 2$; simétricamente por la
izquierda. Los dos estilos importan --- el estructural (construir a
partir de \omterm{def:b1:topology:open}{abiertos} conocidos con uniones e intersecciones finitas)
escala mejor, y el de $\varepsilon$ funciona cuando no se ve ninguna
estructura; y el~\cref{ch:b1:continuity} añadirá un tercero, el más
potente: $U$ es la \omterm{def:b1:logic:map}{imagen recíproca} del \omterm{def:b1:topology:open}{abierto}
$\intoo{2}{+\infty}$ por la función continua $x \mapsto x^2$.
\end{example}

\begin{definition}[Conjunto cerrado]\label{def:b1:topology:closed}
Un \omterm{def:b1:logic:sets}{conjunto} $F \subseteq \R$ es \emph{cerrado}\index{conjunto cerrado}
cuando su complementario $\R \setminus F$ es \omterm{def:b1:topology:open}{abierto}. Por De Morgan y
la~\cref{prop:b1:topology:openstable}: toda intersección de \omterm{def:b1:logic:sets}{conjuntos}
cerrados es cerrada, y las uniones finitas de cerrados son cerradas.
\end{definition}

\begin{theorem}[Caracterización secuencial de los cerrados]\label{thm:b1:topology:seqclosed}
$F$ es \omterm{def:b1:topology:closed}{cerrado} si y solo si, para toda sucesión $(u_n)$ de puntos de
$F$ que converja a algún $\ell \in \R$, el límite $\ell$ pertenece a
$F$. («\omterm{def:b1:topology:closed}{Cerrado}» $=$ «estable por paso al límite».)
\end{theorem}

\begin{proof}
($\Rightarrow$) Sean $F$ \omterm{def:b1:topology:closed}{cerrado}, $u_n \in F$, $u_n \to \ell$, y
supóngase $\ell \notin F$. El complementario es \omterm{def:b1:topology:open}{abierto}: algún
$\intoo{\ell - r}{\ell + r}$ evita $F$. Pero la convergencia mete
$u_n$ en ese \omterm{prop:b1:reals:intervals}{intervalo} para $n$ grande: contradicción con $u_n \in F$.

($\Leftarrow$) Supóngase que $F$ no es \omterm{def:b1:topology:closed}{cerrado}: el complementario no es
\omterm{def:b1:topology:open}{abierto}, luego algún $x \notin F$ no tiene ningún \omterm{prop:b1:reals:intervals}{intervalo}
$\intoo{x - r}{x + r}$ dentro del complementario; tomando
$r = \frac{1}{n+1}$, elíjase $u_n \in F$ con
$\abs{u_n - x} < \frac{1}{n+1}$. Entonces $u_n \in F$ y
$u_n \to x \notin F$: la propiedad secuencial falla.
\end{proof}

\begin{example}[El test secuencial, en los dos sentidos]\label{ex:b1:topology:seqtest}
\emph{\omterm{def:b1:topology:closed}{Cerrado}:} $F = \Z \cup \bigl\{n + \frac1n : n \geq 2\bigr\}$.
Sea $u_k \in F$ con $u_k \to \ell$. La ventana
$\intcc{\ell - 1}{\ell + 1}$ contiene solo finitos puntos de $F$
(finitos enteros, finitos $n + \frac1n$) y, más allá de cierto rango,
todos los $u_k$ están en ella: la sucesión toma entonces finitos
valores y, siendo convergente, acaba siendo constante (como en
el~\cref{exo:b1:topology:3}): $\ell \in F$. \omterm{def:b1:topology:closed}{Cerrado} --- aunque $F$
contenga pares de puntos a distancia $\frac1n$, arbitrariamente
próximos.

\emph{No \omterm{def:b1:topology:closed}{cerrado}:} $G = \bigl\{\frac1m + \frac1n : m, n \in
\N^*\bigr\}$. La sucesión $\frac1n + \frac1n \in G$ tiende a $0$, y
$0 \notin G$ (suma de dos términos positivos): el test secuencial
falla, $G$ no es \omterm{def:b1:topology:closed}{cerrado}. Curiosamente, cada $\frac1m$ \emph{sí}
pertenece a $\overline G \cap G$: en efecto,
$\frac1m = \frac{1}{m+1} + \frac{1}{m(m+1)} \in G$. La idea de cierre:
para demostrar que un \omterm{def:b1:logic:sets}{conjunto} es \omterm{def:b1:topology:closed}{cerrado} hay que controlar
\emph{todas} las sucesiones convergentes a la vez (normalmente con un
argumento de finitud local o de fórmula \omterm{def:b1:topology:closed}{cerrada}); para refutarlo basta
una sucesión bien elegida que se escape --- la asimetría hace fácil la
dirección negativa, y todos los contraejemplos de este capítulo tienen
esa forma de una línea.
\end{example}

\begin{example}\label{ex:b1:topology:closed}
Los segmentos $\intcc{a}{b}$, las semirrectas $\intco{a}{+\infty}$,
los \omterm{def:b1:counting:card}{conjuntos finitos} y $\Z$ (una sucesión convergente de enteros
acaba siendo constante) son \omterm{def:b1:topology:closed}{cerrados}. $\intoc{0}{1}$ no es \omterm{def:b1:topology:open}{abierto}
(falla en $1$) ni \omterm{def:b1:topology:closed}{cerrado} ($\frac 1n \to 0 \notin$ el \omterm{def:b1:logic:sets}{conjunto}): la
mayoría de los \omterm{def:b1:logic:sets}{conjuntos} no son ni una cosa ni la otra. $\R$ y
$\emptyset$ son a la vez \omterm{def:b1:topology:open}{abiertos} y \omterm{def:b1:topology:closed}{cerrados} --- y son los únicos
subconjuntos de $\R$ así (\cref{exo:b1:topology:9}).
\end{example}

\begin{example}[Un abierto ensamblado con infinitas piezas]\label{ex:b1:topology:rminusz}
$\R \setminus \Z = \bigcup_{n \in \Z} \intoo{n}{n+1}$: una unión
infinita de \omterm{prop:b1:reals:intervals}{intervalos} \omterm{def:b1:topology:open}{abiertos}, \omterm{def:b1:topology:open}{abierta} por
la~\cref{prop:b1:topology:openstable} --- de modo que $\Z$ es
\omterm{def:b1:topology:closed}{cerrado} sin necesidad de ningún argumento secuencial. Obsérvese el
reparto de tareas en las reglas de estabilidad: las \emph{uniones} de
\omterm{def:b1:topology:open}{abiertos} pueden ser arbitrarias (cada punto solo necesita su propio
certificado, suministrado por el \omterm{def:b1:logic:sets}{conjunto} que lo contiene), mientras
que las \emph{intersecciones} deben ser finitas (hay que intersecar
los certificados, e infinitos radios pueden encogerse hasta nada).
\Cref{exo:b1:topology:10} mostrará que este ejemplo es la forma
general: todo \omterm{def:b1:topology:open}{abierto} de $\R$ es una unión numerable y disjunta de
\omterm{prop:b1:reals:intervals}{intervalos} \omterm{def:b1:topology:open}{abiertos}.
\end{example}

\section{Interior, clausura, densidad}

\begin{definition}[Interior, clausura, frontera]\label{def:b1:topology:closure}
Sea $A \subseteq \R$.
\begin{itemize}
  \item Un punto $x$ es \emph{interior} a $A$ cuando $A$ es
        \omterm{def:b1:topology:open}{entorno} de $x$; el \emph{interior}
        $\mathring{A}$\index{interior} es el \omterm{def:b1:logic:sets}{conjunto} de los puntos
        interiores.
  \item Un punto $x$ es \emph{adherente} a $A$ cuando todo \omterm{def:b1:topology:open}{entorno}
        de $x$ corta a $A$; la \emph{clausura}
        $\overline{A}$\index{clausura} es el \omterm{def:b1:logic:sets}{conjunto} de los puntos
        adherentes.
  \item La \emph{frontera} es $\partial A = \overline A \setminus
        \mathring A$.\index{frontera}
\end{itemize}
Se tiene $\mathring A \subseteq A \subseteq \overline A$.
\end{definition}

\begin{proposition}[Propiedades principales]\label{prop:b1:topology:closureprops}
\begin{enumerate}
  \item $\mathring A$ es el mayor \omterm{def:b1:topology:open}{abierto} contenido en $A$; $A$ es
        \omterm{def:b1:topology:open}{abierto} si y solo si $A = \mathring A$.
  \item $\overline A$ es el menor \omterm{def:b1:topology:closed}{cerrado} que contiene a $A$; $A$ es
        \omterm{def:b1:topology:closed}{cerrado} si y solo si $A = \overline A$.
  \item (Caracterización secuencial de la adherencia)
        $x \in \overline A$ si y solo si $x$ es límite de una sucesión
        de puntos de $A$.
  \item La complementación intercambia las nociones: $\R \setminus
        \overline A = \bigl(\R \setminus A\bigr)^{\!\circ}$.
\end{enumerate}
\end{proposition}

\begin{proof}
(4) $x \notin \overline A$ $\iff$ algún \omterm{def:b1:topology:open}{entorno} de $x$ evita $A$
$\iff$ algún \omterm{prop:b1:reals:intervals}{intervalo} en torno a $x$ está en $\R \setminus A$ $\iff$
$x$ es \omterm{def:b1:topology:closure}{interior} a $\R \setminus A$.

(1) $\mathring A$ es \omterm{def:b1:topology:open}{abierto}: si $x \in \mathring A$, algún
$\intoo{x-r}{x+r} \subseteq A$; todo punto $y$ de ese \omterm{prop:b1:reals:intervals}{intervalo} tiene
a su alrededor un \omterm{prop:b1:reals:intervals}{intervalo} menor contenido en él y, por tanto, en
$A$: todo el \omterm{prop:b1:reals:intervals}{intervalo} está en $\mathring A$. Todo \omterm{def:b1:topology:open}{abierto}
$U \subseteq A$ consta de puntos \omterm{def:b1:topology:closure}{interiores} de $A$, luego
$U \subseteq \mathring A$: es el mayor. La caracterización de los
\omterm{def:b1:topology:open}{abiertos} se sigue.

(2) Con detalle. Por (4), $\R \setminus \overline A$ es el \omterm{def:b1:topology:closure}{interior}
de $\R \setminus A$, un \omterm{def:b1:topology:open}{abierto} por (1): luego $\overline A$ es
\omterm{def:b1:topology:closed}{cerrado}, y contiene a $A$. Minimalidad: sea $F \supseteq A$
\omterm{def:b1:topology:closed}{cerrado}. Entonces $\R \setminus F$ es \omterm{def:b1:topology:open}{abierto} y está contenido en
$\R \setminus A$, luego, por la maximalidad de (1),
\[
\R \setminus F \subseteq \bigl(\R \setminus A\bigr)^{\!\circ}
= \R \setminus \overline A ,
\]
y, tomando complementarios otra vez, $\overline A \subseteq F$. Así
pues, $\overline A$ es el menor \omterm{def:b1:topology:closed}{cerrado} que lo contiene.
Caracterización: si $A = \overline A$, entonces $A$ es \omterm{def:b1:topology:closed}{cerrado}
(recién visto); y si $A$ es \omterm{def:b1:topology:closed}{cerrado}, es él mismo un \omterm{def:b1:topology:closed}{cerrado} que
contiene a $A$, luego la minimalidad obliga a
$\overline A \subseteq A$, y hay igualdad.

(3) Si $u_n \in A$, $u_n \to x$: todo \omterm{def:b1:topology:open}{entorno} de $x$ contiene algún
$u_n \in A$, luego $x \in \overline A$. Recíprocamente, si
$x \in \overline A$: cada \omterm{prop:b1:reals:intervals}{intervalo} $\intoo{x - \frac{1}{n+1}}{x +
\frac{1}{n+1}}$ corta a $A$ en algún $u_n$, y $u_n \to x$.
\end{proof}

\begin{example}\label{ex:b1:topology:closureexamples}
$\overline{\intoo{0}{1}} = \intcc{0}{1}$;
$\mathring{\intcc{0}{1}} = \intoo{0}{1}$;
$\partial\intoo{0}{1} = \{0, 1\}$. Para $A = \{\frac 1n : n \in \N^*\}$:
$\overline A = A \cup \{0\}$, $\mathring A = \emptyset$,
$\partial A = A \cup \{0\}$. Para $\Q$: por \omterm{def:b1:topology:dense}{densidad}
(\cref{thm:b1:reals:density}), todo real es adherente a $\Q$, luego
$\overline{\Q} = \R$, mientras que $\mathring{\Q} = \emptyset$ (todo
\omterm{prop:b1:reals:intervals}{intervalo} contiene irracionales): la \omterm{def:b1:topology:closure}{frontera} de $\Q$ es todo $\R$.
\end{example}

\begin{example}[Una anatomía completa]\label{ex:b1:topology:anatomy}
Sea $A = \intoc{0}{1} \,\cup\, \bigl(\Q \cap \intoo{2}{3}\bigr)
\,\cup\, \{4\}$. Calculemos los tres \omterm{def:b1:logic:sets}{conjuntos} de
la~\cref{def:b1:topology:closure}, pieza a pieza.

\emph{\omterm{def:b1:topology:closure}{Interior}.} Un punto de $\intoo{0}{1}$ tiene todo un
\omterm{prop:b1:reals:intervals}{intervalo} dentro de $A$: es \omterm{def:b1:topology:closure}{interior}. El punto $1$: todo \omterm{prop:b1:reals:intervals}{intervalo} a
su alrededor se escapa a la derecha de $1$, donde $A$ no tiene nada
hasta $2$: no es \omterm{def:b1:topology:closure}{interior}. Ningún punto de $\Q \cap \intoo{2}{3}$ es
\omterm{def:b1:topology:closure}{interior} (todo \omterm{prop:b1:reals:intervals}{intervalo} contiene irracionales,
\cref{thm:b1:reals:density}); tampoco el aislado $4$. Luego
$\mathring A = \intoo{0}{1}$.

\emph{\omterm{def:b1:topology:closure}{Clausura}.} Límites de puntos de $A$: todo $\intcc{0}{1}$
($0 = \lim \frac1n$ con $\frac 1n \in A$); todo $\intcc{2}{3}$ (todo
real de ahí es límite de racionales del \omterm{prop:b1:reals:intervals}{intervalo}, otra vez la
\omterm{def:b1:topology:dense}{densidad}); y $4$. Nada más: un punto fuera de
$\intcc{0}{1} \cup \intcc{2}{3} \cup \{4\}$ está a distancia positiva
de ese \omterm{def:b1:topology:closed}{cerrado}. Luego $\overline A = \intcc{0}{1} \cup
\intcc{2}{3} \cup \{4\}$.

\emph{\omterm{def:b1:topology:closure}{Frontera}.} $\partial A = \overline A \setminus \mathring A
= \{0, 1\} \cup \intcc{2}{3} \cup \{4\}$.

La idea de cierre: las tres operaciones actúan \emph{localmente} ---
cada pieza de $A$ aporta según su propia naturaleza (un \omterm{prop:b1:reals:intervals}{intervalo}
macizo conserva su \omterm{def:b1:topology:closure}{interior}, una pieza \omterm{def:b1:topology:dense}{densa} pero porosa se convierte
enteramente en \omterm{def:b1:topology:closure}{frontera}, un punto aislado es pura \omterm{def:b1:topology:closure}{frontera}) --- y un
dibujo de $A$ de dos líneas predice todas las respuestas antes de
escribir demostración alguna.
\end{example}

\begin{remark}[Errores frecuentes al razonar con conjuntos de puntos]
(i) \emph{«No \omterm{def:b1:topology:open}{abierto}» no significa «\omterm{def:b1:topology:closed}{cerrado}»}: la mayoría de los
\omterm{def:b1:logic:sets}{conjuntos} no son ninguna de las dos cosas ($\intoc{0}{1}$) y dos son
las dos ($\emptyset$, $\R$) --- \omterm{def:b1:topology:open}{abierto} y \omterm{def:b1:topology:closed}{cerrado} no son opuestos,
sino duales por complementación. (ii) \emph{\omterm{def:b1:topology:closure}{Interior} y \omterm{def:b1:topology:closure}{clausura} no
conmutan}: para $A = \Q$,
\[
\overline{\mathring A} = \overline\emptyset = \emptyset
\qquad\text{mientras que}\qquad
\bigl(\,\overline A\,\bigr)^{\!\circ} = \mathring \R = \R :
\]
los dos operadores iterados difieren tanto como pueden diferir dos
\omterm{def:b1:logic:sets}{conjuntos}. (iii) \emph{Las uniones infinitas de \omterm{def:b1:topology:closed}{cerrados} pueden no
ser \omterm{def:b1:topology:closed}{cerradas}}: $\bigcup_{n\geq1} \intcc{\frac1n}{1} = \intoc{0}{1}$
--- el espejo del contraejemplo de intersecciones de
la~\cref{prop:b1:topology:openstable}. (iv) \emph{\omterm{def:b1:topology:dense}{Denso} no significa
grande}: $\Q$ es \omterm{def:b1:topology:dense}{denso}, numerable y de \omterm{def:b1:topology:closure}{interior} vacío, y su
complementario también es \omterm{def:b1:topology:dense}{denso}; la \omterm{def:b1:topology:dense}{densidad} dice «arbitrariamente
cerca de todo», no «casi todo» --- el \omterm{pb:b1:topology:1}{conjunto de Cantor} del problema
del fin de semana (\cref{pb:b1:topology:1}) señala lo contrario: un
\omterm{def:b1:logic:sets}{conjunto} topológicamente pequeño que es no numerablemente grande.
\end{remark}

\begin{example}[Una clausura calculada exactamente]\label{ex:b1:topology:closurecomputed}
Sea $G = \bigl\{\frac1m + \frac1n : m, n \in \N^*\bigr\}$ (del
\cref{ex:b1:topology:seqtest}). Afirmación:
\[
\overline G = G \,\cup\, \Bigl\{\frac1m : m \in \N^*\Bigr\}
\,\cup\, \{0\} .
\]
($\supseteq$) $\frac1m = \lim_n \bigl(\frac1m + \frac1n\bigr)$ y
$0 = \lim_n \frac2n$: adherentes por la caracterización secuencial.
($\subseteq$) Sea $x = \lim_k \bigl(\frac{1}{m_k} +
\frac{1}{n_k}\bigr)$; ordénese cada par de modo que $m_k \leq n_k$. Si
$(m_k)$ no está acotada, una \omterm{def:b1:seq:subsequence}{subsucesión} tiene $m_k \to \infty$, luego
también $n_k \to \infty$ y $x = 0$. En caso contrario, $(m_k)$ toma
finitos valores, uno de ellos, digamos $m$, infinitas veces; a lo
largo de esa \omterm{def:b1:seq:subsequence}{subsucesión}, $\frac{1}{n_k} \to x - \frac1m$: si $(n_k)$
está acotada, toma algún valor $n$ infinitas veces y
$x = \frac1m + \frac1n \in G$; y si no, $x = \frac1m$. Todos los casos
caen en el \omterm{def:b1:logic:sets}{conjunto} anunciado. La idea de cierre: calcular una
\omterm{def:b1:topology:closure}{clausura} es un análisis por casos al estilo de la compacidad sobre
los índices --- índice acotado significa finitos valores (palomar),
índice no acotado significa que un límite se escapa --- y la respuesta
exhibe la típica estructura en dos capas de los puntos límite: el
\omterm{def:b1:logic:sets}{conjunto}, sus límites de primera generación y el límite de estos, $0$.
\end{example}

\begin{definition}[Densidad, forma topológica]\label{def:b1:topology:dense}
$A$ es \emph{denso}\index{densidad} en $\R$ cuando $\overline A = \R$
--- equivalentemente, todo \omterm{prop:b1:reals:intervals}{intervalo} \omterm{def:b1:topology:open}{abierto} no vacío corta a $A$; y,
equivalentemente (por la~\cref{prop:b1:topology:closureprops}~(3)),
todo real es límite de elementos de $A$. Ejemplos: $\Q$,
$\R \setminus \Q$, los diádicos (\cref{exo:b1:reals:8}), los
\omterm{def:b1:structures:subgroup}{subgrupos} densos (\cref{exo:b1:reals:9}).
\end{definition}

\begin{example}[La densidad es relativa]\label{ex:b1:topology:relativedensity}
«\omterm{def:b1:topology:dense}{Denso}» tal como se ha definido aquí significa \omterm{def:b1:topology:dense}{denso} \emph{en $\R$};
un \omterm{def:b1:logic:sets}{conjunto} puede en cambio ser \omterm{def:b1:topology:dense}{denso} solo en una parte de la recta.
Los diádicos de $\intcc{0}{1}$, es decir, $D \cap \intcc{0}{1}$
(\cref{exo:b1:reals:8}), cortan a todo \omterm{prop:b1:reals:intervals}{intervalo} \omterm{def:b1:topology:open}{abierto} incluido en
$\intcc{0}{1}$, pero por supuesto no tocan $\intoo{2}{3}$: son \omterm{def:b1:topology:dense}{densos}
\emph{en} $\intcc{0}{1}$, es decir, $\overline{D \cap \intcc{0}{1}} =
\intcc{0}{1}$. La frase general «$A$ es \omterm{def:b1:topology:dense}{denso} en $B$» abrevia
$B \subseteq \overline A$ --- nómbrese siempre el \omterm{def:b1:logic:sets}{conjunto} ambiente,
pues los extremos del problema del fin de semana son \omterm{def:b1:topology:dense}{densos} \emph{en
el \omterm{pb:b1:topology:1}{conjunto de Cantor}} y a la vez \omterm{def:b1:topology:dense}{densos} en ninguna parte \emph{de
$\R$}: el mismo \omterm{def:b1:logic:sets}{conjunto}, dos descripciones veraces y de sonido opuesto.
\end{example}

\begin{example}[Manejar la densidad]\label{ex:b1:topology:densitytricks}
Tres movimientos rápidos que se repiten constantemente.
\emph{Agrandar}: si $A$ es \omterm{def:b1:topology:dense}{denso} y $A \subseteq B$, entonces $B$ es
\omterm{def:b1:topology:dense}{denso} (todo \omterm{prop:b1:reals:intervals}{intervalo} ya corta a $A$). \emph{Transportar}: si $A$
es \omterm{def:b1:topology:dense}{denso}, también lo es $\lambda A + \mu$ para $\lambda \neq 0$ ---
un \omterm{prop:b1:reals:intervals}{intervalo} $I$ corta a $\lambda A + \mu$ si y solo si el \omterm{prop:b1:reals:intervals}{intervalo}
$\frac{I - \mu}{\lambda}$ corta a $A$; así, los múltiplos impares de
$10^{-9}$, por ejemplo, son \omterm{def:b1:topology:dense}{densos}. \emph{Intersecar falla}: dos
\omterm{def:b1:logic:sets}{conjuntos} \omterm{def:b1:topology:dense}{densos} pueden no cortarse en absoluto ($\Q$ y
$\R \setminus \Q$): la \omterm{def:b1:topology:dense}{densidad} sobrevive a las uniones y a las
aplicaciones afines, nunca a las intersecciones.
\end{example}

\section{Compacidad de los segmentos}

\begin{theorem}[Los segmentos son secuencialmente compactos]\label{thm:b1:topology:compact}
Sea $a \leq b$. Toda sucesión de puntos de $\intcc{a}{b}$ tiene una
\omterm{def:b1:seq:subsequence}{subsucesión} que converge a un punto \emph{de}
$\intcc{a}{b}$.\index{compacidad}

Más en general, los subconjuntos de $\R$ con esta propiedad (toda
sucesión tiene una \omterm{def:b1:seq:subsequence}{subsucesión} que converge dentro del \omterm{def:b1:logic:sets}{conjunto}) son
exactamente los \omterm{def:b1:logic:sets}{conjuntos} \emph{\omterm{def:b1:topology:closed}{cerrados} y acotados}.
\end{theorem}

\begin{proof}
Una sucesión de $\intcc{a}{b}$ está acotada, luego
Bolzano--Weierstrass (\cref{thm:b1:seq:bw}) extrae una \omterm{def:b1:seq:subsequence}{subsucesión}
convergente; y su límite se queda en $\intcc{a}{b}$ porque los
segmentos son \omterm{def:b1:topology:closed}{cerrados} (\cref{thm:b1:topology:seqclosed}).

Caso general. \emph{(\omterm{def:b1:topology:closed}{Cerrado} y acotado $\Rightarrow$ compacto)}: sean
$F$ \omterm{def:b1:topology:closed}{cerrado} y acotado y $(u_n)$ una sucesión de $F$. La acotación de
$F$ acota la sucesión, luego Bolzano--Weierstrass extrae
$u_{\varphi(n)} \to \ell$; y $\ell \in F$ porque $F$ es \omterm{def:b1:topology:closed}{cerrado} y la
\omterm{def:b1:seq:subsequence}{subsucesión} es una sucesión convergente de puntos de $F$
(\cref{thm:b1:topology:seqclosed}): las dos hipótesis se consumen una
cada una, la acotación para la existencia del límite y el ser \omterm{def:b1:topology:closed}{cerrado}
para su pertenencia. \emph{(Compacto $\Rightarrow$ \omterm{def:b1:topology:closed}{cerrado} y
acotado)}: si $F$ no está acotado, tómese $u_n \in F$ con
$\abs{u_n} \geq n$; toda \omterm{def:b1:seq:subsequence}{subsucesión} no está acotada y, por tanto,
diverge (\cref{prop:b1:seq:first}): no hay ninguna \omterm{def:b1:seq:subsequence}{subsucesión}
convergente. Si $F$ no es \omterm{def:b1:topology:closed}{cerrado}, tómese $u_n \in F$ con
$u_n \to \ell \notin F$ (\cref{thm:b1:topology:seqclosed}): toda
\omterm{def:b1:seq:subsequence}{subsucesión} converge a $\ell \notin F$, luego ninguna converge
\emph{en} $F$.
\end{proof}

\begin{example}[Compactos encajados]\label{ex:b1:topology:nested}
Un primer ejercicio con el teorema. Sean
$K_0 \supseteq K_1 \supseteq K_2 \supseteq \dots$ subconjuntos
compactos (\omterm{def:b1:topology:closed}{cerrados} y acotados) no vacíos de $\R$. Entonces
$\bigcap_n K_n \neq \emptyset$. En efecto, tómese $x_n \in K_n$ para
cada $n$: la sucesión vive en el compacto $K_0$, luego una
\omterm{def:b1:seq:subsequence}{subsucesión} $x_{\varphi(n)}$ converge a cierto $x$
(\cref{thm:b1:topology:compact}). Para cada $m$ fijo, los términos
$x_{\varphi(n)}$ con $\varphi(n) \geq m$ están todos en el \omterm{def:b1:topology:closed}{cerrado}
$K_m$, luego el límite $x$ está en $K_m$
(\cref{thm:b1:topology:seqclosed}); y como $m$ era arbitrario,
$x \in \bigcap_n K_n$. Un compañero útil: si un \omterm{def:b1:topology:open}{abierto} $U$ contiene
$\bigcap_n K_n$, entonces $U \supseteq K_n$ para algún $n$ ---
aplíquese el mismo argumento a puntos $x_n \in K_n \setminus U$; el
límite $x$ estaría en $\bigcap K_n \subseteq U$ y, sin embargo, ser
$U$ \omterm{def:b1:topology:open}{abierto} obliga a que $x_{\varphi(n)} \in U$ a la larga, una
contradicción. Los dos \omterm{def:b1:logic:statement}{enunciados} fallan sin compacidad:
$\bigcap_n \intoo{0}{\frac 1n} = \emptyset$ y
$\bigcap_n \intco{n}{+\infty} = \emptyset$. La idea de cierre: la
compacidad convierte una cadena infinita de afirmaciones de no
vacuidad en un único punto límite --- es la herramienta que sobrevive
al paso a la intersección infinita, y el problema del fin de semana
(\cref{pb:b1:topology:1}) se apoyará en ella dos veces.
\end{example}

\begin{omfigure}
\begin{tikzpicture}[scale=9.5]
  \node[left, font=\small] at (-0.02,0) {$C_0$};
  \fill[omDef] (0,-0.011) rectangle (1,0.011);
  \node[left, font=\small] at (-0.02,-0.09) {$C_1$};
  \fill[omDef] (0,-0.101) rectangle (0.3333,-0.079);
  \fill[omDef] (0.6667,-0.101) rectangle (1,-0.079);
  \node[left, font=\small] at (-0.02,-0.18) {$C_2$};
  \fill[omDef] (0,-0.191) rectangle (0.1111,-0.169);
  \fill[omDef] (0.2222,-0.191) rectangle (0.3333,-0.169);
  \fill[omDef] (0.6667,-0.191) rectangle (0.7778,-0.169);
  \fill[omDef] (0.8889,-0.191) rectangle (1,-0.169);
  \node[left, font=\small] at (-0.02,-0.27) {$C_3$};
  \fill[omDef] (0,-0.281) rectangle (0.0370,-0.259);
  \fill[omDef] (0.0741,-0.281) rectangle (0.1111,-0.259);
  \fill[omDef] (0.2222,-0.281) rectangle (0.2593,-0.259);
  \fill[omDef] (0.2963,-0.281) rectangle (0.3333,-0.259);
  \fill[omDef] (0.6667,-0.281) rectangle (0.7037,-0.259);
  \fill[omDef] (0.7407,-0.281) rectangle (0.7778,-0.259);
  \fill[omDef] (0.8889,-0.281) rectangle (0.9259,-0.259);
  \fill[omDef] (0.9630,-0.281) rectangle (1,-0.259);
  \node[below, font=\small] at (0,-0.30) {$0$};
  \node[below, font=\small] at (0.3333,-0.30) {$\tfrac13$};
  \node[below, font=\small] at (0.6667,-0.30) {$\tfrac23$};
  \node[below, font=\small] at (1,-0.30) {$1$};
  \draw[omThm, thick, -Stealth] (0.5,-0.042) -- (0.5,-0.062);
  \draw[omThm, thick, -Stealth] (0.5,-0.132) -- (0.5,-0.152);
  \draw[omThm, thick, -Stealth] (0.5,-0.222) -- (0.5,-0.242);
\end{tikzpicture}

{\small Las cuatro primeras etapas de la construcción por tercios
centrales: cada segmento de $C_n$ pierde su tercio central \omterm{def:b1:topology:open}{abierto},
dejando los $2^{n+1}$ segmentos de $C_{n+1}$, de longitud total
$(\frac23)^{n+1}$. El \omterm{pb:b1:topology:1}{conjunto de Cantor} $C = \bigcap_n C_n$ --- objeto
del problema del fin de semana \cref{pb:b1:topology:1} --- es el
residuo compacto no vacío que garantiza el argumento de compactos
encajados de más arriba: longitud cero y, aun así, sobreviven no
numerablemente muchos puntos.}
\end{omfigure}

\begin{remark}
Este es el motor que hay detrás del teorema de los valores extremos
(\cref{ch:b1:continuity}) y del teorema de Heine sobre la continuidad
uniforme. El nombre «compacto» adquirirá su definición general (por
recubrimientos) en el segundo año; sobre $\R$, la compacidad
secuencial es todo lo que hace falta, y «compacto $=$ \omterm{def:b1:topology:closed}{cerrado} $+$
acotado» es el \omterm{def:b1:logic:statement}{enunciado} que hay que recordar.
\end{remark}

\begin{example}[Fronteras bajo uniones]\label{ex:b1:topology:boundaryunion}
Siempre $\partial(A \cup B) \subseteq \partial A \cup \partial B$: un
punto de $\partial(A \cup B)$ tiene todo \omterm{def:b1:topology:open}{entorno} cortando a
$A \cup B$ (luego a $A$ o a $B$, y a uno de ellos infinitas veces) y
cortando al complementario de $A \cup B$, que está en los dos
complementarios --- una comprobación breve sitúa entonces el punto en
$\partial A$ o en $\partial B$. La inclusión puede ser
espectacularmente estricta: con $A = \Q$ y $B = \R\setminus\Q$,
\[
\partial(A \cup B) = \partial \R = \emptyset ,
\qquad
\partial A \cup \partial B = \R \cup \R = \R :
\]
dos \omterm{def:b1:logic:sets}{conjuntos} deshilachados pueden pegarse en uno sin costuras,
aniquilándose sus \omterm{def:b1:topology:closure}{fronteras}. La idea de cierre: los \omterm{def:b1:topology:closure}{interiores} y las
\omterm{def:b1:topology:closure}{clausuras} se comportan \emph{monótonamente} bajo uniones e
intersecciones, pero las \omterm{def:b1:topology:closure}{fronteras} no --- trátese $\partial$ como una
cantidad derivada ($\overline A \setminus \mathring A$), nunca como un
operador con álgebra propia.
\end{example}

\begin{remark}[Perspectivas dentro de este volumen]
El vocabulario construido aquí se consume dos veces más en este libro.
En el~\cref{ch:b1:continuity}, todo teorema es un \omterm{def:b1:logic:statement}{enunciado} de
topología disfrazado: el teorema del valor intermedio dice que las
aplicaciones continuas conservan la propiedad de ser \omterm{prop:b1:reals:intervals}{intervalo}, y el
de los valores extremos, que conservan la compacidad --- y las
demostraciones llaman por su nombre a los
\cref{thm:b1:topology:seqclosed,thm:b1:topology:compact}. En
el~\cref{ch:b1:multivar}, las mismas definiciones se releen en $\R^2$
con discos en lugar de \omterm{prop:b1:reals:intervals}{intervalos}: \omterm{def:b1:topology:open}{abiertos}, \omterm{def:b1:topology:closure}{clausuras} y compacidad
se trasladan palabra por palabra, y el teorema de los valores extremos
en dos variables cabalga de nuevo sobre Bolzano--Weierstrass
(extráigase en cada coordenada). La única noción que \emph{no} se
generaliza sin dolor es el propio \omterm{prop:b1:reals:intervals}{intervalo} --- en el plano, la
conexión sustituye a la convexidad, historia que empieza con el «solo
$\emptyset$ y $\R$ son \omterm{def:b1:topology:open}{abiertos} y \omterm{def:b1:topology:closed}{cerrados}»
del~\cref{exo:b1:topology:9}.
\end{remark}

\section{Ejercicios}

\begin{exercise}[$\star$]\label{exo:b1:topology:1}
Dígase, para cada \omterm{def:b1:logic:sets}{conjunto}, si es \omterm{def:b1:topology:open}{abierto}, \omterm{def:b1:topology:closed}{cerrado}, las dos cosas o
ninguna (con justificación): $\intoo{0}{1} \cup \intoo{2}{3}$;
$\;\intco{0}{1}$; $\;\{0\} \cup \intcc{1}{2}$; $\;\R \setminus \Z$;
$\;\Q \cap \intoo{0}{1}$.
\end{exercise}

\begin{exercise}[$\star$]\label{exo:b1:topology:2}
Determínense $\mathring A$, $\overline A$ y $\partial A$ para:
$A = \intoc{0}{1} \cup \{2\}$; $\;A = \R \setminus \Q$;
$\;A = \bigl\{\frac{(-1)^n n}{n+1} : n \in \N\bigr\}$.
\end{exercise}

\begin{exercise}[$\star$]\label{exo:b1:topology:3}
Demuéstrese que un \omterm{def:b1:counting:card}{conjunto finito} es \omterm{def:b1:topology:closed}{cerrado}, primero con
complementarios y después con la caracterización secuencial.
\end{exercise}

\begin{exercise}[$\star$]\label{exo:b1:topology:4}
Demuéstrese que, para todos $A, B \subseteq \R$,
$\overline{A \cup B} = \overline A \cup \overline B$. Véase con un
ejemplo que $\overline{A \cap B}$ puede diferir de
$\overline A \cap \overline B$.
\end{exercise}

\begin{exercise}[$\star\star$]\label{exo:b1:topology:5}
Sea $u_n \to \ell$ en $\R$. Demuéstrese que el \omterm{def:b1:logic:sets}{conjunto}
$\{u_n : n \in \N\} \cup \{\ell\}$ es \omterm{def:b1:topology:closed}{cerrado} (y, por tanto, compacto
si se añade que está acotado --- y lo está).
\end{exercise}

\begin{exercise}[$\star\star$]\label{exo:b1:topology:6}
Sean $U$ \omterm{def:b1:topology:open}{abierto} y $A$ arbitrario. Demuéstrese que $U + A = \{u + a\}$
es \omterm{def:b1:topology:open}{abierto}. Dedúzcase que la suma de un \omterm{def:b1:topology:open}{abierto} y un \omterm{def:b1:logic:sets}{conjunto}
cualquiera es \omterm{def:b1:topology:open}{abierta}, y contrástese: exhíbanse dos \omterm{def:b1:logic:sets}{conjuntos}
\emph{\omterm{def:b1:topology:closed}{cerrados}} cuya suma no sea \omterm{def:b1:topology:closed}{cerrada}. \emph{(Pruébese con $\Z$ y
$\sqrt 2\,\Z$, y el~\cref{exo:b1:reals:9}.)}
\end{exercise}

\begin{exercise}[$\star\star$]\label{exo:b1:topology:7}
Un punto $x \in A$ es \emph{aislado} en $A$ cuando algún \omterm{def:b1:topology:open}{entorno} de
$x$ corta a $A$ solo en $x$. Demuéstrese que todo punto de $\Z$ es
aislado en $\Z$, que $A = \{\frac1n\}$ tiene todos sus puntos aislados
y, aun así, $\overline A \neq A$, y que un \omterm{def:b1:logic:sets}{conjunto} cuyos puntos sean
todos aislados tiene \omterm{def:b1:topology:closure}{interior} vacío.
\end{exercise}

\begin{exercise}[$\star\star$]\label{exo:b1:topology:8}
Demuéstrese que la \omterm{def:b1:topology:closure}{clausura} de un \omterm{def:b1:logic:sets}{conjunto} acotado está acotada y que
$\sup \overline A = \sup A$ para $A$ no vacío y acotado
superiormente. Dedúzcase que $\sup A \in \overline A$: el \omterm{def:b1:reals:bounds}{supremo} es
siempre adherente.
\end{exercise}

\begin{exercise}[$\star\star\star$]\label{exo:b1:topology:9}
Demuéstrese que los únicos subconjuntos de $\R$ que son a la vez
\omterm{def:b1:topology:open}{abiertos} y \omterm{def:b1:topology:closed}{cerrados} son $\emptyset$ y $\R$.
\emph{Indicación: supóngase $A$ \omterm{def:b1:topology:open}{abierto} y \omterm{def:b1:topology:closed}{cerrado}, con
$A \neq \emptyset$ y $\R \setminus A \neq \emptyset$; tómense $a \in
A$, $b \notin A$, digamos $a < b$, y considérese $s = \sup\,(A \cap
\intcc{a}{b})$; decídase si $s$ puede estar en $A$ o en su complementario.}
\end{exercise}

\begin{exercise}[$\star\star\star$]\label{exo:b1:topology:10}
(Estructura de los \omterm{def:b1:topology:open}{abiertos}) Sea $U \subseteq \R$ \omterm{def:b1:topology:open}{abierto} y no vacío.
Para $x \in U$, sea $I_x$ la unión de todos los \omterm{prop:b1:reals:intervals}{intervalos} \omterm{def:b1:topology:open}{abiertos}
que contienen a $x$ y están contenidos en $U$. Demuéstrese que $I_x$
es un \omterm{prop:b1:reals:intervals}{intervalo} \omterm{def:b1:topology:open}{abierto}, que dos \omterm{def:b1:logic:sets}{conjuntos} $I_x$, $I_y$ son iguales o
disjuntos, y que $U$ es unión de \emph{numerablemente muchos}
\omterm{prop:b1:reals:intervals}{intervalos} \omterm{def:b1:topology:open}{abiertos} disjuntos dos a dos \emph{(tómese un racional en cada uno)}.
\end{exercise}

\begin{exercise}[$\star\star$]\label{exo:b1:topology:11}
Un punto $x \in \R$ es \emph{punto de acumulación} de $A$ cuando todo
\omterm{def:b1:topology:open}{entorno} de $x$ corta a $A \setminus \{x\}$; su \omterm{def:b1:logic:sets}{conjunto} es el
\emph{\omterm{def:b1:logic:sets}{conjunto} derivado} $A'$. Demuéstrese que
$\overline A = A \cup A'$ y que $A$ es \omterm{def:b1:topology:closed}{cerrado} si y solo si
$A' \subseteq A$. Determínese $A'$ para $A = \{\frac 1n : n \in \N^*\}$,
para $A = \Z$ y para $A = \Q$.
\end{exercise}

\begin{exercise}[$\star\star\star$]\label{exo:b1:topology:12}
Para $A \subseteq \R$ no vacío, defínase
$d_A(x) = \inf\{\abs{x - a} : a \in A\}$. Demuéstrese:
\begin{enumerate}
  \item $\abs{d_A(x) - d_A(y)} \leq \abs{x - y}$ para todos $x, y$
        ($d_A$ es $1$-lipschitziana);
  \item $d_A(x) = 0$ si y solo si $x \in \overline A$; en particular,
        si $F$ es \omterm{def:b1:topology:closed}{cerrado} y $x \notin F$, entonces $d_F(x) > 0$;
  \item para todo $\varepsilon > 0$, el \omterm{def:b1:logic:sets}{conjunto} $V_\varepsilon =
        \{x : d_F(x) < \varepsilon\}$ es \omterm{def:b1:topology:open}{abierto}, contiene a $F$, y
        $\bigcap_{\varepsilon > 0} V_\varepsilon = F$ para $F$
        \omterm{def:b1:topology:closed}{cerrado}: todo \omterm{def:b1:topology:closed}{cerrado} es una intersección numerable de
        \omterm{def:b1:topology:open}{abiertos}.
\end{enumerate}
\end{exercise}

\section{Problema: el conjunto de Cantor, pequeño y enorme a la vez}

\begin{problem}[{Problema del fin de semana --- el \omterm{pb:b1:topology:1}{conjunto de Cantor}
de los tercios centrales: longitud cero, no numerable, perfecto y
$C + C = \intcc{0}{2}$}]
\label{pb:b1:topology:1}
Quítese a $\intcc{0}{1}$ su tercio central \omterm{def:b1:topology:open}{abierto}, después el tercio
central de cada segmento restante, y repítase para siempre: lo que
sobrevive es el \emph{\omterm{pb:b1:topology:1}{conjunto de Cantor}} $C$, la fábrica fundamental
de contraejemplos del análisis. Este problema lo construye, lo lee con
la maquinaria de la base $3$ del~\cref{pb:b1:reals:1} y establece su
retrato paradójico: longitud total cero y, aun así, no numerable;
\omterm{def:b1:topology:closure}{interior} vacío y, aun así, sin ningún punto aislado; totalmente
disconexo y, aun así, $C + C$ llena todo el segmento
$\intcc{0}{2}$.\index{conjunto de Cantor} Formalmente:
$C_0 = \intcc{0}{1}$, y $C_{n+1}$ se obtiene de $C_n$ suprimiendo el
tercio central \omterm{def:b1:topology:open}{abierto} de cada segmento de $C_n$; finalmente,
$C = \bigcap_{n \geq 0} C_n$. En todo el problema, un \emph{código
ternario} de $x \in \intcc{0}{1}$ es cualquier cadena de cifras
$(d_k)_{k\geq1}$ con $d_k \in \{0, 1, 2\}$ cuyo valor
$\sup_n \sum_{k=1}^n d_k 3^{-k}$ sea igual a $x$ --- se admiten los
códigos impropios (finalmente $2$); por el~\cref{pb:b1:reals:1}
(preguntas 9--11), todo $x \in \intcc{0}{1}$ tiene uno o dos códigos,
y dos exactamente cuando $x = m/3^N \in \intoo{0}{1}$.

\medskip
\noindent\textbf{Parte I --- La construcción.}
\begin{enumerate}
  \item Descríbanse $C_1$ y $C_2$ explícitamente como uniones de
        segmentos, y demuéstrese por inducción que $C_n$ es unión
        disjunta de $2^n$ segmentos \omterm{def:b1:topology:closed}{cerrados}, cada uno de longitud
        $3^{-n}$.
  \item Véase que $C$ es \omterm{def:b1:topology:closed}{cerrado} y acotado ---y, por tanto, compacto
        (\cref{thm:b1:topology:compact})--- y no vacío, y que todo
        extremo de todo segmento de todo $C_n$ pertenece a $C$.
  \item La longitud total de $C_n$ es $\bigl(\frac23\bigr)^n$.
        Dedúzcase que, para todo $\varepsilon > 0$, el \omterm{def:b1:logic:sets}{conjunto} $C$ se
        puede recubrir con finitos segmentos de longitud total
        $\leq \varepsilon$: el \omterm{pb:b1:topology:1}{conjunto de Cantor} tiene
        \emph{longitud cero}.
  \item Véase que un \omterm{prop:b1:reals:intervals}{intervalo} contenido en $C$ tiene longitud
        $\leq 3^{-n}$ para todo $n$ y, por tanto, es un \omterm{def:b1:logic:sets}{conjunto}
        unitario o vacío: $\mathring C = \emptyset$. Siendo \omterm{def:b1:topology:closed}{cerrado}
        con \omterm{def:b1:topology:closure}{interior} vacío, $C$ es \emph{\omterm{def:b1:topology:dense}{denso} en ninguna parte}.
\end{enumerate}

\medskip
\noindent\textbf{Parte II --- El código ternario.}
\begin{enumerate}[resume]
  \item Demuéstrese la recursión de autosemejanza
        \[
        C_{n+1} = \tfrac13 C_n \,\cup\,
        \bigl(\tfrac23 + \tfrac13 C_n\bigr),
        \qquad\text{luego}\qquad
        C = \tfrac13 C \,\cup\, \bigl(\tfrac23 + \tfrac13
        C\bigr),
        \]
        siendo las dos piezas disjuntas: $C$ son dos copias de sí
        mismo a escala $\frac13$.
  \item Demuéstrese por inducción sobre $n$: $x \in C_n$ si y solo si
        $x$ tiene un código ternario cuyas primeras $n$ cifras están
        en $\{0, 2\}$. Dedúzcase, usando que $x$ tiene a lo sumo dos
        códigos: $x \in C$ si y solo si $x$ tiene un código
        \emph{sin ninguna cifra igual a $1$} (un código \emph{libre
        de $1$}).
  \item Códigos en acción: dense códigos libres de $1$ para $0$, $1$,
        $\frac13$ y $\frac23$; véase que
        $\frac14 = (0.\overline{02})_3$ y
        $\frac34 = (0.\overline{20})_3$, de modo que los dos están en
        $C$; y compruébese que $\frac14$ \emph{no} es extremo de
        ningún $C_n$ (los extremos tienen la forma $m/3^n$).
  \item Véase que cada $x \in C$ tiene \emph{exactamente un} código
        libre de $1$ \emph{(cuando $x$ tiene dos códigos, demuéstrese
        que exactamente uno de los dos contiene la cifra $1$)}.
        Conclúyase: la \omterm{def:b1:logic:map}{aplicación} valor es una \omterm{def:b1:logic:inj}{biyección} de las
        cadenas de $\{0,2\}$ sobre $C$.
  \item (Diagonal) Sea $k \mapsto x_k$ una \omterm{def:b1:logic:map}{aplicación} cualquiera
        $\N^* \to C$. Constrúyase una cadena de $\{0, 2\}$ que difiera
        en el índice $k$ del código de $x_k$, y conclúyase que $C$ no
        es numerable --- mientras que, por contraste, la pregunta 3
        dice que es métricamente despreciable.
\end{enumerate}

\medskip
\noindent\textbf{Parte III --- Retrato topológico.}
\begin{enumerate}[resume]
  \item Reúnase lo obtenido hasta aquí: $C$ es compacto, no
        numerable, de longitud cero y \omterm{def:b1:topology:dense}{denso} en ninguna parte. ¿Qué
        única contención $C \subseteq C_n$ sostiene cada propiedad?
  \item ($C$ es perfecto) Sea $x \in C$ de código libre de $1$
        $(d_k)$. Cambiar la cifra $d_n$ ($0 \leftrightarrow 2$)
        produce $x_n \in C$ con $\abs{x_n - x} = 2 \cdot 3^{-n}$.
        Conclúyase que $C$ no tiene ningún punto aislado: todo punto
        de $C$ es límite de \emph{otros} puntos de $C$.
  \item Véase que los extremos de la pregunta 2 forman un subconjunto
        numerable y \omterm{def:b1:topology:dense}{denso} de $C$ \emph{(trúnquese el código tras $n$
        cifras y continúese con $0$; la numerabilidad, como en
        el~\cref{exo:b1:topology:10})}. Conclúyase: el punto típico de
        $C$ --- como $\frac14$ --- \emph{no} es un extremo: los
        extremos son un esqueleto numerable dentro de un cuerpo no
        numerable.
  \item (Totalmente disconexo) Sean $x < y$ en $C$. Elíjase $n$ con
        $3^{-n} < y - x$ y prodúzcase un punto $z \in \intoo{x}{y}$
        con $z \notin C$. Conclúyase que los únicos \omterm{prop:b1:reals:intervals}{intervalos} no
        vacíos contenidos en $C$ son los \omterm{def:b1:logic:sets}{conjuntos} unitarios.
\end{enumerate}

\medskip
\noindent\textbf{Parte IV --- Aritmética de $C$.}
\begin{enumerate}[resume]
  \item Véase que $1 - C = C$ \emph{(¿qué le hace $x \mapsto 1 - x$ a
        un código libre de $1$? recuérdese $1 = (0.\overline{2})_3$)}.
  \item (Suma de códigos) Véase que si $x$, $x'$ tienen códigos
        $(a_k)$, $(b_k)$, entonces $x + x' = \lim_n\,(t_n + t'_n)$,
        donde $t_n, t'_n$ son las sumas parciales. Dedúzcase: todo
        $y \in \intcc{0}{1}$ es el \emph{punto medio} de dos puntos de
        $C$ --- dado un código $(e_k)$ de $y$, elíjanse cifras
        $a_k, b_k \in \{0, 2\}$ con $\frac{a_k + b_k}{2} = e_k$.
  \item Conclúyase $C + C = \intcc{0}{2}$ y, con la pregunta 14,
        $C - C = \intcc{-1}{1}$. Instancia concreta: escríbase $1$
        como suma de los dos puntos no extremos hallados en la
        pregunta 7.
  \item Reflexión: un \omterm{def:b1:logic:sets}{conjunto} de longitud cero cuyo \omterm{def:b1:logic:sets}{conjunto} de
        diferencias llena $\intcc{-1}{1}$. ¿Por qué no hay
        contradicción entre «$C$ es métricamente despreciable» y
        «$C + C$ tiene longitud completa»? (Una frase; piénsese en
        qué controla la longitud y qué no.)
\end{enumerate}

\medskip
\noindent\textbf{Parte V --- Miembros racionales e irracionales.}
\begin{enumerate}[resume]
  \item Combínese la pregunta 8 con el criterio de periodicidad
        del~\cref{pb:b1:reals:1}: un punto de $C$ es racional si y
        solo si su código libre de $1$ es finalmente periódico.
        Ejecútese la división larga en base $3$ para comprobar que
        $\frac1{13} = (0.\overline{002})_3 \in C$.
  \item Prodúzcase un miembro explícitamente \emph{irracional} de $C$:
        el valor del código con $d_k = 2$ en las posiciones
        triangulares $k = \frac{j(j+1)}{2}$ y $d_k = 0$ en las demás.
        Justifíquese la irracionalidad con el argumento de los huecos
        crecientes del~\cref{pb:b1:reals:1} (pregunta 20).
  \item (Sobre un segmento completo) Considérese $h$, que envía el
        punto de $C$ de código libre de $1$ $(d_k)$ al valor de la
        cadena \emph{binaria} $\bigl(\frac{d_k}2\bigr)$, es decir,
        $h(x) = \sup_n \sum_{k=1}^n \frac{d_k}{2}\,2^{-k}$. Véase que
        $h$ envía $C$ \emph{sobre} $\intcc{0}{1}$. Así pues, el
        despreciable $C$ se aplica sobreyectivamente sobre un segmento
        de longitud completa --- una segunda demostración de que $C$
        no es numerable.
  \item (Longitud autosemejante) Supóngase que estuviera definida
        alguna noción de longitud $L$ para $C$ y sus copias
        encogidas, respetando el escalado ($L(\lambda A) = \lambda
        L(A)$), la invariancia por traslación y la aditividad sobre la
        descomposición disjunta de la pregunta 5. Véase que entonces
        $L(C) = \frac23\,L(C)$, lo que obliga a $L(C) = 0$: la
        autosemejanza sola ya condena a $C$ a longitud cero.
\end{enumerate}

\medskip
\noindent\textbf{Parte VI --- Un \omterm{def:b1:arith:prime}{primo} gordo, y la moraleja.}
\begin{enumerate}[resume]
  \item (\omterm{pb:b1:topology:1}{Conjunto de Cantor} gordo) Repítase la construcción, pero
        en la etapa $n$ ($n = 0, 1, 2, \dots$) quítese a cada uno de
        los $2^n$ segmentos actuales solo un \omterm{prop:b1:reals:intervals}{intervalo} central
        \omterm{def:b1:topology:open}{abierto} de longitud $4^{-(n+1)}$. Véase que las longitudes
        $l_n$ de los segmentos cumplen
        $l_{n+1} = \frac{l_n - 4^{-(n+1)}}{2}$,
        $l_n = \frac{2^n + 1}{2\cdot 4^n} > 0$, que el $K = \bigcap
        K_n$ resultante es compacto con \omterm{def:b1:topology:closure}{interior} vacío y que la
        longitud total suprimida es $\sum_{n\geq0} 2^n 4^{-(n+1)} =
        \frac12$. Admitiendo la aditividad (intuitiva, del tercer año)
        de la longitud para uniones finitas de \omterm{prop:b1:reals:intervals}{intervalos}, y usando
        los dos \omterm{def:b1:logic:statement}{enunciados} del~\cref{ex:b1:topology:nested}, véase que
        toda familia finita de \omterm{prop:b1:reals:intervals}{intervalos} \omterm{def:b1:topology:open}{abiertos} que recubra $K$
        tiene longitud total $\geq \frac12$: $K$ es \omterm{def:b1:topology:dense}{denso} en ninguna
        parte pero \emph{no} despreciable. La pequeñez tiene varios
        significados no equivalentes.
  \item (Distancias) Véase que, para $F \subseteq \R$ \omterm{def:b1:topology:closed}{cerrado} y no
        vacío y $x \in \R$, el \omterm{def:b1:reals:bounds}{ínfimo} $d(x, F)$ se \emph{alcanza}
        \emph{(sucesión minimizante más Bolzano--Weierstrass)}.
        Calcúlese después
        \[
        \max_{y \in \intcc{0}{1}} d(y, C) = \frac16 ,
        \]
        alcanzado exactamente en el centro $y = \frac12$ \emph{(un
        punto de un hueco creado en la etapa $n$ dista menos de
        $\frac{3^{-n}}{2}$ de los extremos del hueco, que están en
        $C$)}.
  \item (Todo punto, límite de una \omterm{def:b1:seq:subsequence}{subsucesión}) Usando las
        preguntas 12 y 2, prodúzcase una única sucesión de $C$ cuyo
        \omterm{def:b1:logic:sets}{conjunto} de límites de subsucesiones sea \emph{todo} $C$.
        (Compárese: para una sucesión convergente ese \omterm{def:b1:logic:sets}{conjunto} es un
        solo punto --- $C$ realiza el extremo opuesto entre los
        compactos.)
  \item Síntesis, una frase para cada punto: (i) ¿qué teoremas de este
        capítulo ha consumido realmente la construcción (estabilidad
        de los \omterm{def:b1:topology:closed}{cerrados}, compacidad, caracterizaciones
        secuenciales)? (ii) enumérense las cuatro parejas paradójicas
        del retrato (longitud cero/no numerable, \omterm{def:b1:topology:closed}{cerrado}/\omterm{def:b1:topology:closure}{interior}
        vacío, perfecto/totalmente disconexo, despreciable/$C+C$
        completo); (iii) ¿dónde reaparece $C$ más adelante (la
        escalera del diablo construida sobre $h$ en la teoría de la
        continuidad, y la teoría de la medida del volumen del tercer
        año, donde $C$ separa «numerable» de «despreciable»)?
\end{enumerate}
\end{problem}
